\documentclass[a4paper]{article}

%------------------------------------------------
% الحزم
%------------------------------------------------

\usepackage{enumitem}
\usepackage[margin=2.2cm]{geometry}
\usepackage{amsmath,amssymb}
\usepackage{graphicx}
\usepackage{array}
\usepackage{fancyhdr}
\usepackage[table]{xcolor}
\usepackage{tikz}
\usetikzlibrary{positioning,arrows.meta,shapes.geometric,calc}
\usepackage[skins]{tcolorbox}
\usepackage{polyglossia}

%------------------------------------------------
% اللغات والخطوط
%------------------------------------------------

\setdefaultlanguage[
    calendar=gregorian,
    numerals=maghrib
]{arabic}

\setotherlanguage{english}

% إذا لم تكن الخطوط مثبتة في الجهاز، تُؤخذ من المجلد fonts/
\IfFontExistsTF{Amiri}{%
    \newfontfamily\arabicfont[Script=Arabic, Scale=1.15]{Amiri}
}{%
    \newfontfamily\arabicfont[
        Script=Arabic, Scale=1.15,
        Path=fonts/, Extension=.ttf,
        UprightFont=*-Regular, BoldFont=*-Bold,
        ItalicFont=*-Italic, BoldItalicFont=*-BoldItalic
    ]{Amiri}
}

\IfFontExistsTF{AlBattar}{%
    \newfontfamily\arabicfontsf[Script=Arabic, Scale=1.2]{AlBattar}
    % خط العناوين (يعمل حتى بعد استعمال \textenglish)
    \newfontfamily\khatt[Script=Arabic, Scale=1.2]{AlBattar}
}{%
    \newfontfamily\arabicfontsf[
        Script=Arabic, Scale=1.2, Path=fonts/,
        BoldFont=ae_AlBattar.ttf
    ]{ae_AlBattar.ttf}
    % خط العناوين (يعمل حتى بعد استعمال \textenglish)
    \newfontfamily\khatt[
        Script=Arabic, Scale=1.2, Path=fonts/,
        BoldFont=ae_AlBattar.ttf
    ]{ae_AlBattar.ttf}
}

%------------------------------------------------
% الألوان: لون رئيسي هادئ + لون تنبيه فقط
%------------------------------------------------

\definecolor{asasi}{HTML}{3D6A8E}   % اللون الرئيسي: للرسوم والمخططات
\definecolor{tanbih}{HTML}{C0603A}  % لون التنبيه: للتحذيرات والحالات السلبية
\definecolor{ramadi}{HTML}{5F6B73}  % رمادي للنصوص الثانوية والأسهم
\definecolor{dakin}{HTML}{263238}   % لون العناوين

% أسماء قديمة تُحوَّل إلى اللونين فقط
\colorlet{azraq}{asasi}
\colorlet{fayruzi}{asasi}
\colorlet{banafsaji}{asasi}
\colorlet{akhdar}{asasi}
\colorlet{ahmar}{tanbih}
\colorlet{burtuqali}{asasi}

%------------------------------------------------
% الصناديق
%------------------------------------------------

\tcbset{
    asl/.style={
        enhanced,
        arc=2mm,
        boxrule=0.6pt,
        colback=white,
        colframe=black!18,
        colbacktitle=black!5,
        coltitle=dakin,
        titlerule=0pt,
        left=8pt, right=8pt, top=5pt, bottom=5pt,
        fonttitle=\khatt\bfseries,
        before skip=8pt, after skip=8pt
    }
}

% أهداف الدرس
\newtcolorbox{ahdaf}{asl, title={أهداف الدرس}}

% تعريف (خط رئيسي رفيع على الجانب)
\newtcolorbox{taarif}[1][تعريف]{asl, title={#1},
    borderline east={2pt}{0pt}{asasi}}

% مثال
\newtcolorbox{mithal}[1][مثال]{asl, title={#1}}

% تنبيه (اللون الوحيد البارز في النص)
\newtcolorbox{intabih}[1][انتبه!]{asl, colback=tanbih!5, colframe=tanbih,
    colbacktitle=tanbih, coltitle=white, title={#1}}

% ركن أستاذ اللغة العربية
\newtcolorbox{arabiyya}[1][في قسم اللغة العربية]{asl, title={#1}}

% نص أمر يكتب للأداة
\newtcolorbox{amr}[1][ramadi]{enhanced, arc=1mm, boxrule=0pt,
    colback=black!4, borderline east={2.5pt}{0pt}{#1},
    left=8pt, right=10pt, top=4pt, bottom=4pt,
    before skip=6pt, after skip=6pt}

% الخلاصة
\newtcolorbox{khulasa}{asl, title={خلاصة الدرس}}

% عنوان القسم
\newcounter{qism}
\newcommand{\qism}[1]{%
    \stepcounter{qism}%
    \par\bigskip
    {\large\khatt\bfseries{\color{asasi}\arabic{qism}.}\ {\color{dakin}#1}}\par
    \vspace{-2pt}{\color{black!20}\rule{\linewidth}{0.6pt}}\par\nopagebreak[4]
    \smallskip}

% عنوان فرعي
\newcommand{\far}[1]{\par\smallskip
    {\color{dakin}\khatt\bfseries {\color{asasi}$\blacktriangleleft$}\ #1}\par\nopagebreak[4]}

% كلمة مهمة (خط عريض فقط)
\newcommand{\kalima}[2][]{\textbf{#2}}

%------------------------------------------------
% أنماط الرسوم
%------------------------------------------------

\newcommand{\nodecolor}{}

\tikzset{
    % كتابة النص العربي من اليمين إلى اليسار داخل العقد
    % (لون النص يُعطى بالمفتاح lawn=... وليس text=...)
    every node/.append style={
        execute at begin node={\nodecolor\everypar{\beginR}\beginR},
        execute at end node=\endR},
    lawn/.code={\def\nodecolor{\color{#1}}},
    khatwa/.style={draw=#1, fill=#1!12, very thick, rounded corners=6pt,
        minimum height=1.2cm, text width=2.3cm, align=center,
        font=\khatt\bfseries},
    khatwa/.default=azraq,
    sahm/.style={-{Stealth[length=3mm]}, very thick, color=ramadi},
    raqm/.style={circle, fill=#1, lawn=white, inner sep=2pt,
        font=\small\bfseries},
}

%------------------------------------------------
% إعداد الصفحة
%------------------------------------------------

\renewcommand{\footrulewidth}{1pt}
\renewcommand{\headrulewidth}{1pt}

\rhead{\khatt المدرسة العليا للأساتذة ببوزريعة}
\lhead{\khatt ملحقة بومرداس}

\cfoot{\thepage}
\lfoot{\khatt الأستاذ: د. بن عبد الرزاق يوسف}

\parindent 0pt
\setlength{\headsep}{15pt}
\setlength{\parskip}{4pt}

\pagestyle{fancy}

\setlist[enumerate,1]{label=\arabic*), itemsep=2pt, topsep=3pt}
\setlist[itemize]{itemsep=2pt, topsep=3pt}

% الجداول
\arrayrulecolor{black!22}
\renewcommand{\arraystretch}{1.35}

% رأس جدول ملون
\newcommand{\ras}[1]{\cellcolor{black!7}\color{dakin}\bfseries #1}

% عنوان الدرس
\newcommand{\unwan}[3]{%
    \begin{center}
        {\color{asasi}\rule{\linewidth}{1.2pt}}\par\medskip
        {\khatt\color{ramadi} #1}\par\smallskip
        {\Large\khatt\bfseries\color{dakin} #2}\par\smallskip
        {\khatt\color{ramadi} #3}\par\medskip
        {\color{asasi}\rule{\linewidth}{1.2pt}}
    \end{center}}

%================================================
% بداية الوثيقة
%================================================

\begin{document}

%------------------------------------------------
% العنوان
%------------------------------------------------

\unwan{درس الفصل الرابع}
{الأخلاقيات في الذكاء الاصطناعي}
{المسؤولية ــ الشفافية ــ الإفصاح ــ الأمانة العلمية}

\begin{ahdaf}
في نهاية هذا الدرس، تستطيع أن:
\begin{itemize}
    \item تعرّف المسؤولية والشفافية والإفصاح.
    \item تميّز بين الاستعمال المقبول والاستعمال المرفوض.
    \item تعرّف الانتحال العلمي والانتحال بالذكاء الاصطناعي.
    \item تكتب جملة إفصاح بسيطة في عملك.
\end{itemize}
\end{ahdaf}

%================================================
\qism{لماذا نتحدث عن الأخلاقيات؟}
%================================================

الذكاء الاصطناعي أداة قوية وسريعة.
لكن الأداة لا تقرر وحدها كيف نستعملها.
\kalima{الإنسان} هو الذي يختار: يستعملها بأمانة، أو يستعملها للغش.

وأستاذ المستقبل قدوة لتلاميذه، فسلوكه مع هذه الأدوات درس في حد ذاته.

%================================================
\qism{ثلاثة مبادئ أساسية}
%================================================

\begin{center}
\begin{tikzpicture}[
    amud/.style={draw=#1, fill=#1!12, very thick,
        minimum width=3.6cm, minimum height=2.6cm,
        text width=3.3cm, align=center}]
    % السقف
    \fill[banafsaji!80] (-6.6,2.1) -- (6.6,2.1) -- (0,3.4) -- cycle;
    \node[lawn=white, font=\khatt\bfseries] at (0,2.55) {الاستعمال الأخلاقي};
    % الأعمدة
    \node[amud=azraq, lawn=azraq!70!black]   at (4.4,0.6)
        {{\khatt\bfseries المسؤولية}\\[3pt]\small أنا مسؤول عمّا أقدّمه};
    \node[amud=akhdar, lawn=akhdar!70!black]  at (0,0.6)
        {{\khatt\bfseries الشفافية}\\[3pt]\small أوضّح كيف استعملت الأداة};
    \node[amud=burtuqali, lawn=burtuqali!70!black] at (-4.4,0.6)
        {{\khatt\bfseries الإفصاح}\\[3pt]\small أصرّح باستعمالها};
    % القاعدة
    \fill[ramadi!70] (-6.6,-0.85) rectangle (6.6,-1.3);
\end{tikzpicture}
\end{center}

\begin{taarif}
\textbf{المسؤولية:} المستخدم مسؤول عن استعمال الأداة، وعن كل معلومة يعتمد عليها.

\textbf{الشفافية:} توضيح طريقة استعمال الذكاء الاصطناعي عند الحاجة.

\textbf{الإفصاح:} التصريح باستعمال أداة الذكاء الاصطناعي عندما يكون ذلك مطلوبًا أو مناسبًا.
\end{taarif}

\begin{intabih}
الاستعانة بالذكاء الاصطناعي \textbf{لا تنقل المسؤولية} إلى الأداة.
إذا قدّمت معلومة خاطئة، فأنت المسؤول عنها، لا الأداة.
\end{intabih}

%================================================
\qism{الانتحال العلمي}
%================================================

\begin{taarif}
\textbf{الانتحال العلمي:} أن ينسب الشخص إلى نفسه عملًا أو أفكارًا ليست له،
دون ذكر مصدرها.

\textbf{الانتحال بالذكاء الاصطناعي} \textenglish{(AI-Plagiarism)}:
أن ينسخ الطالب نصًا أنتجته الأداة، ويقدمه على أنه عمله الشخصي.
\end{taarif}

\begin{mithal}[حالة]
نسخ طالب بحثًا كاملًا من أداة ذكاء اصطناعي، ثم كتب اسمه عليه وقدّمه للأستاذ.
هذا \kalima[ahmar]{انتحال}: العمل ليس عمله، ولم يتحقق من صحته، ولم يصرّح بمصدره.
\end{mithal}

%================================================
\qism{مقياس الاستعمال: من المقبول إلى المرفوض}
%================================================

\begin{center}
\begin{tikzpicture}
    \shade[left color=ahmar!80, right color=akhdar!80, rounded corners=4pt]
        (-7.5,0) rectangle (7.5,0.6);
    \node[font=\khatt\bfseries, lawn=akhdar] at (6,1)  {مقبول};
    \node[font=\khatt\bfseries, lawn=ramadi] at (0,1) {مقبول بشروط};
    \node[font=\khatt\bfseries, lawn=ahmar] at (-6,1) {مرفوض};
    \node[text width=4.2cm, align=center, font=\small] at (5.3,-0.8)
        {أفهم مفهومًا صعبًا};
    \node[text width=4.2cm, align=center, font=\small] at (0,-0.8)
        {أنشئ مسودة ثم أراجعها وأصرّح بذلك};
    \node[text width=4.2cm, align=center, font=\small] at (-5.3,-0.8)
        {أنسخ إجابة كاملة وأقدّمها باسمي};
    \foreach \x/\c in {5.3/akhdar, 0/burtuqali, -5.3/ahmar}
        \fill[\c] (\x,0.3) circle (4pt);
\end{tikzpicture}
\end{center}

\begin{center}
\begin{tabular}{>{\raggedleft\arraybackslash}m{9.2cm}
                >{\centering\arraybackslash}m{4.5cm}}
\hline
\ras{الاستعمال} & \ras{الموقف الأخلاقي} \\
\hline
استعمال الأداة لفهم قاعدة نحوية صعبة.
    & \kalima[akhdar]{مقبول} \\ \hline
طلب أفكار من الأداة، ثم كتابة الإجابة بنفسي.
    & \kalima[akhdar]{مقبول} \\ \hline
إنشاء مسودة بالأداة، ثم مراجعتها والتحقق منها والإفصاح.
    & \kalima[burtuqali]{مقبول بشروط} \\ \hline
الاعتماد على معلومات الأداة دون التحقق منها.
    & \kalima[ahmar]{غير مسؤول} \\ \hline
نسخ إجابة كاملة وتقديمها باسمي.
    & \kalima[ahmar]{انتحال مرفوض} \\ \hline
\end{tabular}
\end{center}

%================================================
\qism{كيف أكتب جملة إفصاح؟}
%================================================

جملة الإفصاح قصيرة وصادقة. نضعها في آخر البحث أو في مقدمته.

\begin{amr}[akhdar]
\textbf{نموذج:} استعملتُ أداة \textenglish{ChatGPT} لتنظيم أفكار الفصل الثاني
واقتراح خطة أولية.
راجعتُ كل المعلومات وتحققت منها في المصادر المذكورة، والصياغة النهائية من عملي.
\end{amr}

\begin{arabiyya}
عندما تصبح أستاذًا، علّم تلاميذك هذه القاعدة البسيطة:
«الأداة تساعدني على التفكير، ولا تفكر بدلًا مني».
\end{arabiyya}

%================================================
\qism{قواعد ذهبية للاستعمال الأخلاقي}
%================================================

\begin{center}
\begin{tikzpicture}[
    qaida/.style={draw=#1, fill=#1!10, thick, rounded corners=5pt,
        text width=12.5cm, align=right, minimum height=0.9cm, inner sep=6pt}]
    \node[qaida=azraq]     (q1) at (0,0)    {أتحقق من كل معلومة مهمة قبل استعمالها.};
    \node[qaida=akhdar]    (q2) at (0,-1.1) {أكتب العمل النهائي بأسلوبي وأفكاري.};
    \node[qaida=burtuqali] (q3) at (0,-2.2) {أصرّح باستعمال الأداة عندما يكون ذلك مطلوبًا.};
    \node[qaida=banafsaji] (q4) at (0,-3.3) {أحترم قوانين المؤسسة وتعليمات الأستاذ.};
    \foreach \i/\c in {1/azraq,2/akhdar,3/burtuqali,4/banafsaji}
        \node[raqm=\c, minimum size=0.7cm] at ($(q\i.east)+(0.55,0)$) {\i};
\end{tikzpicture}
\end{center}

\begin{khulasa}
\begin{itemize}
    \item المبادئ الثلاثة: \textbf{المسؤولية}، \textbf{الشفافية}، \textbf{الإفصاح}.
    \item الأداة لا تتحمل المسؤولية؛ \textbf{المستخدم} هو المسؤول.
    \item نسخ عمل الأداة وتقديمه باسمي \textbf{انتحال}.
    \item الاستعمال الجيد: أفهم، أتحقق، أكتب بنفسي، وأصرّح.
\end{itemize}
\end{khulasa}

\end{document}
